\chapter*{Preface}
In $\Z[\sqrt{-5}]$ the number $6$ is both $2\cdot3$ and $(1+\sqrt{-5})(1-\sqrt{-5})$, and unique factorisation is lost.
Kummer and Dedekind found it again by factoring \emph{ideals} instead of numbers, and with that repair a whole world
opened. This book is a journey through that world: the arithmetic of number fields, of the $p$-adic numbers, and of
curves, up to the modular forms that tie them together.

Volume~12 ended with a question. Its last chapters set the zeros of the Riemann zeta function beside the eigenvalues of
large random matrices, and the spacings, the moments and even the maximum matched. Yet $\zeta$ is not random: it is one
function, with an Euler product and a functional equation. Is it alone? This volume answers that it is one of a large
family. Number fields have their Dedekind zeta functions, characters have their $L$-functions, curves over finite fields
have zeta functions for which the Riemann hypothesis is a theorem, and modular forms have $L$-functions with their own
functional equations. Each member is built here from the arithmetic it counts, and each functional equation that
earlier volumes admitted is proved here or marked as a sketch.

The first part opens with quadratic fields. The ring of integers is computed by the trace and the norm, and the ideals
repair $6=2\cdot3=(1+\sqrt{-5})(1-\sqrt{-5})$ in $\Z[\sqrt{-5}]$. The same part then leaves the quadratic shortcut.
Integers of a number field form a lattice whose volume, the discriminant, decides which primes
ramify. Every prime number splits into prime ideals in a way that a single polynomial modulo $p$ reveals, and in Galois
extensions this splitting is governed by one group element, the Frobenius. Minkowski's convex body theorem shows that
class groups are finite and that unit groups are as large as they can be, and Kummer's cyclotomic fields give the
first serious attack on Fermat's Last Theorem. The same chapter proves Fermat's Last Theorem for polynomials in a few lines,
by Mason's theorem, and asks which integers are sums of three cubes, a question with a congruence answer, a heuristic,
and a computer search that settled the last case below $100$ only in 2019.

The second and third parts look at one prime at a time, then at all of them together. They begin by constructing $\Q_p$: an infinite string of residues modulo $p,p^2,p^3,\dots$, with Ostrowski's list of absolute values and Hensel's lemma in Newton's form. Extensions of those fields are
classified by two numbers, $e$ and $f$. The local--global principle of Hasse and Minkowski decides every quadratic
equation by congruences and inequalities, and settles which integers are sums of three squares. Zeta functions count
ideals, compute class numbers, and prove Chebotarev's theorem that every conjugacy class of a Galois group occurs as a
Frobenius with the frequency its size predicts, measured in Dirichlet density. Ad\`eles put the real and $p$-adic worlds on an equal footing, and
Tate's thesis derives the functional equations of the Riemann zeta function and of Dirichlet $L$-functions from Fourier analysis on the ad\`eles of $\Q$. The functional equation of $\zeta_K$ for a general number field is only sketched. Class field theory
then describes all abelian extensions of $\Q$: the Kronecker--Weber theorem is proved in full, and the two main theorems
of the general theory, Artin reciprocity and the existence theorem, and with them the Hilbert class field, are stated,
clearly marked as admitted, and used to decide which primes are of the form $x^2+ny^2$.

The fourth part is arithmetic geometry. Plane curves meet as Bézout predicts, and Riemann--Roch measures their functions.
On an elliptic curve the points form a group, finitely generated over $\Q$ by Mordell's theorem, proved here for every curve, and over a
finite field their number obeys Hasse's bound, a special case of the Riemann hypothesis for curves that we prove with
Bombieri and Stepanov's elementary method. The book ends with modular forms: the discriminant function
$\Delta=q\prod(1-q^n)^{24}$, its coefficients $\tau(n)$ (Ramanujan's function, not the divisor function of Volume~6, Chapter~2), Hecke operators and the multiplicativity they explain,
Hecke's proof that their $L$-functions have a functional equation, and a form whose coefficients count the points of an
elliptic curve modulo every prime.

The aim throughout is to understand why the theorems are true, and to see them at work. Every result that is stated as a theorem is proved, apart from Artin reciprocity, the existence theorem and the Hilbert class field, which are marked admitted in their headings; the results that rest on them say so. Several further results appear only in remarks, marked quoted, and are not proved: Mazur's torsion theorem, the Eichler--Shimura theorem and the modularity theorem, the Sato--Tate theorem, Weil's theorem for a general curve, Chebotarev's theorem in natural density, and the second case of Fermat's Last Theorem for regular primes. The functional equation of $\zeta_K$ for a general number field is a sketch. Pause
and think boxes ask you to guess before reading on. Each chapter ends with exercises, graded from warm-ups
($\star$) to small gems ($\star\star\star$), each with a hint and a complete solution that first explains how to find
the idea and then writes it up.

The beauty of the subject is in its coincidences that turn out not to be coincidences. Whether a prime $p\equiv1\pmod 3$ splits completely in
$\Q(\sqrt[3]2)$ is decided by writing $p=x^2+27y^2$. The number of points of $y^2=x^3-x$ modulo $p$ is predicted by
Gaussian integers. The coefficients of a power series in $q$ count the points of a curve. Each of these is a shadow of
one large structure, and the chapters that follow try to make that structure visible.

The series asks whether pure order and pure chance can be told apart, and why the primes look random yet obey laws.
This volume supplies the order: class numbers, the class number formula, Chebotarev's theorem and the functional
equations of $L$-functions. Volume~14 uses them to measure how far the primes can be predicted, and Volume~15 asks how
class groups and elliptic curves behave on average, when the field or the curve is chosen at random.
